% Extracted from p09-mechanistic-loops
\subsection{What Do Extra Loops Compute?}
\label{ch-mech}

\paragraph{Introduction}

The loop sweep is the most informative cheap experiment you can run on
a recurrent language model. For our 44M-parameter OpenMythos model
(43,974,722 parameters; dim 512; 6 prelude + 2 coda layers; 4 loops;
8 experts; GQA; trained on TinyStories for 120M tokens, loss
9.5242$\rightarrow$1.5263), the sweep over 256 held-out sequences
gives (full results in Section~\ref{sec:loopsweep}, Table~\ref{tab:sweep}):

Two facts stand out. First, the gains are enormous and concentrated:
one extra loop (1$\rightarrow$2) more than halves perplexity; the
trained depth reaches 4.41. The iterations are doing substantial,
non-redundant work. Second, the dynamics saturate hard: doubling past
the trained depth slightly hurts, then nothing changes at all through
16 loops---yet generations stay grammatical and coherent throughout,
so this is not collapse, it is convergence.

The natural next question---\emph{what is each loop computing?}---has
barely been asked for recurrent language models. Feedforward
interpretability has tools for this: induction heads
\citep{elhage2021mathematical} localise in-context learning to
specific heads; the Universal Transformer literature
\citep{dehghani2019universal} speculated about iterative refinement
but rarely measured it per-iteration. We propose to turn the probe
toolkit onto the loop dimension itself.

\paragraph{Hypotheses}

We state three testable hypotheses, each tied to the observed sweep.

\paragraph{H1: Coarse-to-fine refinement.} Early loops resolve local
choices (lexicon, agreement, short-range syntax); later loops handle
long-range coherence (entity consistency, discourse structure). Under
H1, the 29.38$\rightarrow$11.18 drop (loop 1$\rightarrow$2) should be
dominated by local-likelihood improvements, while the
11.18$\rightarrow$4.41 drop (loop 2$\rightarrow$4) should show gains
concentrated on tokens whose predictability depends on distant
context. This is testable by stratifying per-token loss improvements
by dependency length---a measurement, not a speculation.

\paragraph{H2: Fixed-point convergence.} Beyond the trained depth, the
recurrent block converges to an approximate fixed point:
$h_{k} \approx h_{k-1}$ for $k > 4$, which explains both the plateau
(4.90$\rightarrow$4.91$\rightarrow$4.91) and the preserved coherence
(the fixed point is a sensible representation, just not a better
one). Under H2, representation similarity between consecutive loops
should approach 1.0 for $k \geq 8$ while remaining well below 1.0 for
$k \leq 4$.

\paragraph{H3: Progressive circuit formation.} Attention circuits that
implement in-context behaviours---induction heads in the sense of
\citet{elhage2021mathematical}---form progressively across loops
rather than appearing fully formed. Under H3, induction-head scores
(prefix-matching behaviour in attention patterns) should rise
monotonically over loops 1--4 and then stabilise, mirroring the
perplexity curve.

Note the tension H1 and H2 jointly create, which is exactly what makes
this interesting: the same block must implement \emph{changing}
computation early (refinement) and \emph{stationary} computation late
(convergence). Characterising the transition is the core empirical
target.

\paragraph{Proposed Probe Protocol}

All experiments below are proposed. The model, data (256 held-out
TinyStories sequences of length 256), and loop settings
$\{1,2,4,8,12,16\}$ are fixed from the companion study; only the
measurements are new.

\paragraph{P1: Per-loop representation similarity.} For each sequence,
record the recurrent block's output $h_k$ at each loop $k$. Compute
centred kernel alignment (CKA) and mean cosine similarity between
$(h_k, h_{k+1})$ for $k = 1 \ldots 15$ at fixed total depth 16.
\emph{Predicted:} similarity well below 1.0 for $k < 4$, rising
sharply, and $\approx$1.0 for $k \geq 8$ (H2). A control---similarity
between $h_k$ at matched loop index across different total depths
(e.g.\ loop 2 of a 4-loop run vs loop 2 of a 16-loop run)---tests
whether loop-indexed computation is depth-invariant.

\paragraph{P2: Attention-pattern evolution.} For each loop, compute
induction-head scores per head following \citet{elhage2021mathematical}
(attention paid to the token following a previous occurrence of the
current token). Track the distribution of scores across loops.
\emph{Predicted:} median induction score rises over loops 1--4 and
flattens thereafter (H3); heads that score highly should also show
the largest cross-loop attention-map change early and the smallest
late.

\paragraph{P3: Per-loop early-exit classifiers.} Freeze the model;
train a linear probe on each $h_k$ to predict the next token, and
record probe perplexity per loop. Additionally, train probes for
auxiliary targets: part-of-speech of the next token (local syntax)
and the identity of the story's protagonist (global coherence).
\emph{Predicted:} next-token probe perplexity tracks the observed
29.38$\rightarrow$4.41 curve; the syntax probe saturates by loop 2
while the protagonist probe keeps improving to loop 4 (H1).

\paragraph{P4: Loop interventions.} At fixed total depth 16: (a) skip
loop $k$ (replace $h_k$ with $h_{k-1}$) for each $k$ and measure the
perplexity delta---this prices each iteration's marginal contribution;
(b) freeze attention to its loop-4 pattern for all $k > 4$ and test
whether the 4.90 plateau persists, distinguishing representational
from attentional convergence.

\paragraph{What Would Change Our Minds}

The protocol is designed to be falsifiable. If P1 shows high
similarity even between loops 1 and 2, H2's convergence story is
wrong and something else---perhaps the coda, not the recurrent
block---explains the gains. If P3's probes are flat across loops
while output perplexity still drops 29.38$\rightarrow$4.41, the
refinement is happening outside the probed representations (prelude
or coda), which would redirect the investigation. If induction scores
(P2) are already maximal at loop 1, then in-context circuitry is
established immediately and later loops do something orthogonal to
the induction-head picture. Each outcome narrows the hypothesis
space; that is the point.

\paragraph{Discussion}

Why does this matter beyond one 44M model? If recurrent refinement
is real and characterisable---if we can say ``loop 2 fixes syntax,
loop 4 fixes discourse''---then loop count becomes an
\emph{interpretable} compute knob rather than a blind
hyperparameter. It would also inform distillation into recurrent
students (companion paper on loop-aware distillation, this series):
per-loop supervision targets should match what each loop is actually
for. Conversely, if the probes show the block converging to a fixed
point with no structured per-loop division of labour, that reframes
recurrence as learned depth with shared parameters---valuable, but
not ``refinement'' in any mechanistic sense, and the test-time
extrapolation story needs rethinking.

\paragraph{Limitations}

This paper proposes; it does not report. The hypotheses are grounded
in one model's behavioural signature on one dataset (TinyStories),
and synthetic children's stories may not exercise the circuitry
(e.g.\ induction heads) that matters at scale. Probes are
correlational: rising probe accuracy does not prove the model
\emph{uses} the probed information. The companion behavioural numbers
are real; everything about mechanism here is conjecture awaiting the
protocol.

\paragraph{Conclusion}

Perplexity 29.38 $\rightarrow$ 4.41 $\rightarrow$ 4.91 across loops is
a behavioural fact in search of a mechanistic explanation. We have
posed the explanation as three falsifiable hypotheses and specified
the four-probe protocol that would test them. The loops are doing
something; it is time to look at what.
